%% kapa: exam=2 round=21 year=2010 period=1 subject=appraisal_practice subject_ko=감정평가실무 seq=3 number=3 points=15
%% source: Q-Net 감정평가사 기출문제 (cst00309) artlSeq=5024183 file=extracted/1008/5024183/2010년 제21회 감정평가사 2차시험 기출문제/1교시.hwp
%% conversion: kordoc 4.24.51 -> tools/gen2.py (auto); review=proofread-sonnet; flags=none
\begin{longquestion}{【문제 3】}
총 3인의 조합원으로 구성된 재개발조합은 분양계획의 조정을 위하여 아래와 같은 자료를 수집하였다. 구체적으로 조합원 부담의 감소를 위하여 분양계획 1안을 2안으로 변경하고자 한다. 이때 각 1안 및 2안의 비례율을 산정하여, 2안으로 변경할 경우 유리한 조합원의 순서를 판별하되, 산출과정 및 그 이유를 설명하시오. (15점) \par
<자료 1> 조합원별 종전자산 평가액 \par
1) 김한국 : 80,000,000원 \par
2) 이대한 : 140,000,000원 \par
3) 박조선 : 180,000,000원 \par
<자료 2> 건축계획 및 사업관련비용 등 \par
1) 단위세대 당 105m\textsuperscript{2} 면적의 총 10세대를 건축할 예정이며, 이중 3세대는 조합원 분양분이고, 7세대는 일반분양 예정임. \par
{\emergencystretch=1em
2) 본 사업의 진행을 위해서는 기존주택 등의 철거비로 100,000,000원이 소요될 전망이고, 신축공사비로 1,400,000,000원의 지출이 예상됨. \par}
<자료 3> 분양예정가격 \par
1) 1안 : 조합원에 대한 분양예정가격은 단위세대 당 160,000,000원, 일반분양예정가격은 단위세대 당 200,000,000원임. \par
2) 2안 : 조합원에 대한 분양예정가격은 단위세대 당 140,000,000원, 일반분양예정가격은 단위세대 당 200,000,000원임. \par
\end{longquestion}
